% Generated draft: requires mathematical and source review.
\documentclass[11pt]{article}
\usepackage[margin=25mm]{geometry}
\usepackage{fontspec}
\setmainfont{Latin Modern Roman}
\usepackage{amsmath,amssymb,mathtools,unicode-math}
\setmathfont{Latin Modern Math}
\usepackage{enumitem,needspace,xurl,hyperref}
\setlength{\parindent}{0pt}
\setlength{\parskip}{5pt}
\setlength{\emergencystretch}{3em}
\newcommand{\N}{\mathbb{N}}
\newcommand{\Z}{\mathbb{Z}}
\newcommand{\Q}{\mathbb{Q}}
\newcommand{\R}{\mathbb{R}}
\newcommand{\C}{\mathbb{C}}
\newcommand{\F}{\mathbb{F}}
\newcommand{\A}{\mathbb{A}}
\newcommand{\PP}{\mathbb{P}}
\newcommand{\E}{\mathbb{E}}
\newcommand{\eps}{\varepsilon}
\newcommand{\dd}{\,\mathrm{d}}
\begin{document}
\section*{UnsolvedMath ID: 30001436}
\textbf{Source problem code:} catalogue entry
\subsection{Definitions and mathematical statement}
Let \(G\) be a finite group, with identity element \(1\). Its order \(|G|\) is odd if \(2\nmid |G|\). Define its derived series by
\[
G^{(0)}=G,\qquad G^{(n+1)}=[G^{(n)},G^{(n)}]\quad(n\geq 0),
\]
where \([H,H]\) is the subgroup generated by the commutators \(x^{-1}y^{-1}xy\) with \(x,y\in H\). The group \(G\) is solvable if \(G^{(r)}=\{1\}\) for some integer \(r\geq 0\).

The supplied normalized question is whether
\[
\forall\,G\text{ finite},\qquad 2\nmid |G|\ \Longrightarrow\ \exists r\in\Z_{\geq 0}\text{ such that }G^{(r)}=\{1\}.
\]
This includes the trivial group.

A simple group is a nontrivial group whose only normal subgroups are itself and the trivial subgroup. A composition series of a finite group is a chain
\[
\{1\}=G_0\triangleleft G_1\triangleleft\cdots\triangleleft G_m=G
\]
whose successive quotients \(G_i/G_{i-1}\) are simple; these quotients are its composition factors. The solvability question can equivalently be expressed by requiring every composition factor of every finite odd-order group to be cyclic of odd prime order, or by asking whether there are no finite \emph{nonabelian} simple groups of odd order.

\textbf{Unresolved source qualification.} The supplied record instead calls the absence of \emph{nontrivial} finite simple groups of odd order equivalent to solvability, and its original wording asks for the absence of nontrivial simple composition factors. These formulations are not equivalent under the definitions above: a cyclic group of odd prime order is nontrivial and simple, and is its own composition factor. Thus the displayed solvability question preserves the supplied normalized target, while the nonabelian-simple formulation states its mathematical equivalence explicitly. Whether the source intended ``nonabelian'' where the record says ``nontrivial'' requires review; that intention is not assumed here.
\subsection{Short English statement}
The normalized question asks whether every finite group with an odd number of elements is solvable. Its equivalent simple-group formulation excludes nonabelian simple groups of odd order. The record instead excludes all nontrivial simple groups, which also excludes cyclic groups of odd prime order and therefore needs correction or clarification. The supplied historical context and open-status label require review.
\subsection{Sources}
\begin{itemize}
\item Supplied local UnsolvedMath registry record, permanent ID \(30001436\), title \emph{The Odd-Order Simple-Group Problem}, source problem code \texttt{catalogue entry}; collection: \emph{Oberwolfach Reports Open Problems}. The original, clean, and normalized statement fields are the evidence for this draft.
\item Supplied source citation: \emph{Disciplines and Styles in Pure Mathematics, 1800--2000} (2010), \emph{Oberwolfach Reports}, pp.\ 581--665, DOI 10.4171/owr/2010/12. Supplied source URL and literature-source link, labeled ``Oberwolfach Report 12/2010'': \url{https://doi.org/10.4171/owr/2010/12}. No exact passage or problem page is supplied.
\item Supplied literature assessment, checked \textbf{2026-08-21}: the extracted sentence is described as a historical reference to an open problem identified by Burnside, without enough context to identify a resolved theorem. The supplied status is ``open.'' This is supplied dated curation, not an independently verified literature assessment; the background explicitly requests independent verification of time-sensitive status and cited claims.
\item The record supplies \texttt{statement\_status}: \texttt{corrected\_verified} and reports that the mathematical content was checked against the official report and normalized without changing the source-level claim. This is supplied verification metadata, not verification performed for this draft. No proposer is named; the listed proposed year is 2010, also the report year, and does not establish the historical date of origin.
\end{itemize}
\end{document}
